% !TeX program = XeLaTeX
% !TeX root = ../../shloka.tex
% --meta--
% deity: Shiva
% type: Ashtakam
% chandas: Vasantatilaka
% wiki_title: Vishvanatha Ashtakam
% --end-meta--
% Text: श्रीगङ्गाकल्पोक्तपूजाविधिः (Idkidu, 2015), pp. 49-50 (verse order as printed).

\sect{विश्वनाथाष्टकम्}

\fourlineindentedshloka
{गङ्गातरङ्गकमनीयजटाकलापं}
{गौरीनिरन्तरविभूषितवामभागम्}
{नारायणप्रियमनङ्गमदापहारं}
{वाराणसीपुरपतिं भज विश्वनाथम्} %॥१॥

\fourlineindentedshloka
{वाचामगोचरमनेकगुणस्वरूपं}
{वागीशविष्णुसुरसेवितपादपीठम्}
{वामेन विग्रहवरेण कलत्रवन्तं}
{वाराणसीपुरपतिं भज विश्वनाथम्} %॥२॥

\fourlineindentedshloka
{भूताधिपं भुजगभूषणभूषिताङ्गं}
{व्याघ्राजिनाम्बरधरं जटिलं त्रिनेत्रम्}
{पाशाङ्कुशाभयवरप्रदशूलपाणिं}
{वाराणसीपुरपतिं भज विश्वनाथम्} %॥३॥

\fourlineindentedshloka
{शीतांशुशोभितकिरीटविराजमानं}
{भालेक्षणानलविशोषितपञ्चबाणम्}
{नागाधिपारचितभासुरकर्णपूरं}
{वाराणसीपुरपतिं भज विश्वनाथम्} %॥४॥

\fourlineindentedshloka
{तेजोमयं सकलनिष्कलमद्वितीयम्}
{आनन्दकन्दमपराजितमप्रमेयम्}
{नादात्मकं सगुणनिर्गुणमात्मरूपं}
{वाराणसीपुरपतिं भज विश्वनाथम्} %॥५॥

\fourlineindentedshloka
{पञ्चाननं दुरितमत्तमतङ्गजानां}
{नागान्तकं दनुजपुङ्गवपन्नगानाम्}
{दावानलं मरणशोकजराटवीनां}
{वाराणसीपुरपतिं भज विश्वनाथम्} %॥६॥

\fourlineindentedshloka
{आशां विहाय परिहृत्य परस्य निन्दां}
{पापे रतिं च सुनिवार्य मनः समाधौ}
{आदाय हृत्कमलमध्यगतं परेशं}
{वाराणसीपुरपतिं भज विश्वनाथम्} %॥७॥

\fourlineindentedshloka
{रागादिदोषरहितं स्वजनानुरागं}
{वैराग्यशान्तिनिलयं गिरिजासहायम्}
{माधुर्यधैर्यसुभगं गरलाभिरामं}
{वाराणसीपुरपतिं भज विश्वनाथम्} %॥८॥

\fourlineindentedshloka
{वाराणसीपुरपतेः स्तवनं शिवस्य}
{व्याख्यातमष्टकमिदं पठते मनुष्यः}
{विद्यां श्रियं विपुलसौख्यमनन्तकीर्तिं}
{सम्प्राप्य देहविलये लभते च मोक्षम्} %॥९॥

\twolineshloka
{विश्वनाथाष्टकमिदं यः पठेच्छिवसन्निधौ}
{शिवलोकमवाप्नोति शिवेन सह मोदते} %॥१०॥

{॥इति श्री-विश्वनाथाष्टकं सम्पूर्णम्॥}
